
\documentclass{article}
\usepackage{colortbl}
\usepackage{makecell}
\usepackage{multirow}
\usepackage{supertabular}

\begin{document}

\newcounter{utterance}

\centering \large Interaction Transcript for game `cladder', experiment `full\_v1.5\_default', episode 5168 with refineonly{-}Qwen3.5{-}27B.
\vspace{24pt}

{ \footnotesize  \setcounter{utterance}{1}
\setlength{\tabcolsep}{0pt}
\begin{supertabular}{c@{$\;$}|p{.15\linewidth}@{}p{.15\linewidth}p{.15\linewidth}p{.15\linewidth}p{.15\linewidth}p{.15\linewidth}}
    \# & \multicolumn{2}{c}{Player} && \multicolumn{2}{c}{Game Master} \\
    \hline

    \theutterance \stepcounter{utterance}  
    & & & \multicolumn{4}{p{0.6\linewidth}}{
        \cellcolor[rgb]{0.9,0.9,0.9}{
            \makecell[{{p{\linewidth}}}]{
                \texttt{\tiny{[P1$\langle$GM]}}
                \texttt{You are an expert in causal inference. The following question is not a typical commonsense query, but rather a meticulously designed question created by a professor specializing in causal inference, intended to assess the students' mastery of the course content.} \\
\texttt{Imagine a self{-}contained, hypothetical world with only the following conditions, and without any unmentioned factors or causal relationships: Talent has a direct effect on elite institution admission status. Brown eyes has a direct effect on elite institution admission status. The overall probability of talent is 97\%. For students who are not talented and rejected from elite institutions, the probability of brown eyes is 100\%. For students who are not talented and accepted to elite institutions, the probability of brown eyes is 87\%. For students who are talented and rejected from elite institutions, the probability of brown eyes is 97\%. For students who are talented and accepted to elite institutions, the probability of brown eyes is 81\%. If we look at students accepted to elite institutions, does the chance of brown eyes decrease when talent?} \\
\texttt{Start your answer with 'Yes' or 'No', followed by additional reasoning or evidence to support your explanation.} \\
            }
        }
    }
     \\ \\

    \theutterance \stepcounter{utterance}  
    & & \multicolumn{4}{p{0.6\linewidth}}{
        \cellcolor[rgb]{0.95,0.95,0.95}{
            \makecell[{{p{\linewidth}}}]{
                \texttt{\tiny{[P1$\leftrightarrow$P]}}
                \texttt{{[}DRAFT{]}} \\
\texttt{Yes} \\
\\ 
\texttt{**Reasoning:**} \\
\\ 
\texttt{To determine if the chance of brown eyes decreases when talent is present among students accepted to elite institutions, we must compare the conditional probability of having brown eyes given acceptance and talent (\$P(\textbackslash{}text\{Brown\} \textbackslash{}mid \textbackslash{}text\{Accepted\}, \textbackslash{}text\{Talent\})\$) against the conditional probability of having brown eyes given acceptance and no talent (\$P(\textbackslash{}text\{Brown\} \textbackslash{}mid \textbackslash{}text\{Accepted\}, \textbackslash{}neg\textbackslash{}text\{Talent\})\$).} \\
\\ 
\texttt{The problem explicitly provides these conditional probabilities:} \\
\texttt{1.  **For students who are not talented and accepted:** The probability of brown eyes is **87\%**.} \\
\texttt{    \$\$P(\textbackslash{}text\{Brown\} \textbackslash{}mid \textbackslash{}text\{Accepted\}, \textbackslash{}neg\textbackslash{}text\{Talent\}) = 0.87\$\$} \\
\texttt{2.  **For students who are talented and accepted:** The probability of brown eyes is **81\%**.} \\
\texttt{    \$\$P(\textbackslash{}text\{Brown\} \textbackslash{}mid \textbackslash{}text\{Accepted\}, \textbackslash{}text\{Talent\}) = 0.81\$\$} \\
\\ 
\texttt{Comparing the two values:} \\
\texttt{\$\$0.81 < 0.87\$\$} \\
\\ 
\texttt{Since the probability of having brown eyes is lower for the talented group (81\%) than for the non{-}talented group (87\%) within the subset of accepted students, the chance of brown eyes indeed decreases when talent is present in this specific context.} \\
\\ 
\texttt{*Note on Causal Structure:* This scenario illustrates a classic case of **collider bias** (or Berkson's paradox). "Elite institution admission" acts as a collider because it is caused by both "Talent" and "Brown eyes." Conditioning on the collider (looking only at accepted students) induces a spurious negative association between the two causes, even if they are independent in the general population. The data provided confirms this induced negative correlation: among the accepted, being talented is associated with a lower likelihood of having brown eyes.} \\
            }
        }
    }
    & \\ \\

    \theutterance \stepcounter{utterance}  
    & \multicolumn{4}{p{0.6\linewidth}}{
        \cellcolor[rgb]{0.561,0.737,0.561}{
            \makecell[{{p{\linewidth}}}]{
                \texttt{\tiny{[P1$\rangle$GM]}}
                \texttt{Yes} \\
\\ 
\texttt{**Reasoning:**} \\
\\ 
\texttt{To determine if the chance of brown eyes decreases when talent is present among students accepted to elite institutions, we must compare the conditional probability of having brown eyes given acceptance and talent (\$P(\textbackslash{}text\{Brown\} \textbackslash{}mid \textbackslash{}text\{Accepted\}, \textbackslash{}text\{Talent\})\$) against the conditional probability of having brown eyes given acceptance and no talent (\$P(\textbackslash{}text\{Brown\} \textbackslash{}mid \textbackslash{}text\{Accepted\}, \textbackslash{}neg\textbackslash{}text\{Talent\})\$).} \\
\\ 
\texttt{The problem explicitly provides these conditional probabilities:} \\
\texttt{1.  **For students who are not talented and accepted:** The probability of brown eyes is **87\%**.} \\
\texttt{    \$\$P(\textbackslash{}text\{Brown\} \textbackslash{}mid \textbackslash{}text\{Accepted\}, \textbackslash{}neg\textbackslash{}text\{Talent\}) = 0.87\$\$} \\
\texttt{2.  **For students who are talented and accepted:** The probability of brown eyes is **81\%**.} \\
\texttt{    \$\$P(\textbackslash{}text\{Brown\} \textbackslash{}mid \textbackslash{}text\{Accepted\}, \textbackslash{}text\{Talent\}) = 0.81\$\$} \\
\\ 
\texttt{Comparing the two values:} \\
\texttt{\$\$0.81 < 0.87\$\$} \\
\\ 
\texttt{Since the probability of having brown eyes is lower for the talented group (81\%) than for the non{-}talented group (87\%) within the subset of accepted students, the chance of brown eyes indeed decreases when talent is present in this specific context.} \\
\\ 
\texttt{*Note on Causal Structure:* This scenario illustrates a classic case of **collider bias** (or Berkson's paradox). "Elite institution admission" acts as a collider because it is caused by both "Talent" and "Brown eyes." Conditioning on the collider (looking only at accepted students) induces a spurious negative association between the two causes, even if they are independent in the general population. The data provided confirms this induced negative correlation: among the accepted, being talented is associated with a lower likelihood of having brown eyes.} \\
            }
        }
    }
    & & \\ \\

    \theutterance \stepcounter{utterance}  
    & & & \multicolumn{2}{p{0.3\linewidth}}{
        \cellcolor[rgb]{0.9,0.9,0.9}{
            \makecell[{{p{\linewidth}}}]{
                \texttt{\tiny{[GM$|$GM]}}
                \texttt{Yes} \\
            }
        }
    }
    & & \\ \\

    \theutterance \stepcounter{utterance}  
    & & & \multicolumn{2}{p{0.3\linewidth}}{
        \cellcolor[rgb]{0.9,0.9,0.9}{
            \makecell[{{p{\linewidth}}}]{
                \texttt{\tiny{[GM$|$GM]}}
                \texttt{yes} \\
            }
        }
    }
    & & \\ \\

    \theutterance \stepcounter{utterance}  
    & & & \multicolumn{2}{p{0.3\linewidth}}{
        \cellcolor[rgb]{0.9,0.9,0.9}{
            \makecell[{{p{\linewidth}}}]{
                \texttt{\tiny{[GM$|$GM]}}
                \texttt{game\_result = WIN} \\
            }
        }
    }
    & & \\ \\

\end{supertabular}
}

\end{document}
